\hypertarget{extreme-low-latency-hardware-mastery}{%
\chapter{EXTREME LOW LATENCY \& HARDWARE
MASTERY}\label{extreme-low-latency-hardware-mastery}}

To achieve sub-microsecond latency, you must program the hardware, not
just the language.

\hypertarget{cpu-architecture-cache-topology}{%
\subsection{31.1 CPU Architecture \& Cache
Topology}\label{cpu-architecture-cache-topology}}

\begin{itemize}
\tightlist
\item
  \textbf{L1 Cache}: \textasciitilde32KB, 3-4 cycles. Per core.
\item
  \textbf{L2 Cache}: \textasciitilde256KB-1MB, 10-12 cycles. Per core.
\item
  \textbf{L3 Cache}: \textasciitilde10MB+, 40-70 cycles. Shared across
  cores.
\item
  \textbf{RAM}: 100+ cycles.
\end{itemize}

\textbf{Optimization Goal}: Stay in L1/L2. \textbf{Technique}: Minimize
object size, use contiguous memory (arrays), align data to cache lines
(64 bytes).

\hypertarget{numa-non-uniform-memory-access}{%
\subsection{31.2 NUMA (Non-Uniform Memory
Access)}\label{numa-non-uniform-memory-access}}

On multi-socket servers, accessing RAM attached to another CPU socket is
slow. * \textbf{Solution}: Pin threads to cores. Allocate memory on the
local node. * \textbf{Tool}:
\passthrough{\lstinline!numactl --cpunodebind=0 --membind=0 ./app!}

\hypertarget{compiler-optimizations-the-free-lunch}{%
\subsection{31.3 Compiler Optimizations (The ``Free
Lunch'')}\label{compiler-optimizations-the-free-lunch}}

\begin{itemize}
\tightlist
\item
  \passthrough{\lstinline!-O3!}: Aggressive optimization.
\item
  \passthrough{\lstinline!-march=native!}: Use instructions available on
  the build machine (AVX2, AVX-512).
\item
  \passthrough{\lstinline!-flto!} (Link Time Optimization): Optimize
  across translation units (inlining across .cpp files).
\item
  \textbf{PGO (Profile Guided Optimization)}:

  \begin{enumerate}
  \def\labelenumi{\arabic{enumi}.}
  \tightlist
  \item
    Compile with \passthrough{\lstinline!-fprofile-generate!}.
  \item
    Run the app (training run).
  \item
    Recompile with \passthrough{\lstinline!-fprofile-use!}.
  \end{enumerate}
\end{itemize}

\hypertarget{lock-free-stack-implementation-wait-free-push}{%
\subsection{31.4 Lock-Free Stack Implementation (Wait-Free
Push)}\label{lock-free-stack-implementation-wait-free-push}}

A classic interview and system component.

\begin{lstlisting}[language={C++}]
template<typename T>
struct Node {
    T data;
    Node* next;
    Node(const T& d) : data(d), next(nullptr) {}
};

template<typename T>
class LockFreeStack {
    std::atomic<Node<T>*> head{nullptr};

public:
    void push(const T& data) {
        Node<T>* new_node = new Node<T>(data);
        new_node->next = head.load(std::memory_order_relaxed);
        
        // CAS Loop
        while (!head.compare_exchange_weak(
            new_node->next, 
            new_node,
            std::memory_order_release, 
            std::memory_order_relaxed));
    }

    bool pop(T& result) {
        Node<T>* old_head = head.load(std::memory_order_acquire);
        
        while (old_head && !head.compare_exchange_weak(
            old_head,
            old_head->next,
            std::memory_order_acquire,
            std::memory_order_relaxed));
            
        if (!old_head) return false;
        
        result = old_head->data;
        // Note: Deletion in lock-free requires Hazard Pointers or RCU!
        // Leaking here for simplicity of example.
        return true;
    }
};
\end{lstlisting}

\hypertarget{measurable-performance-targets}{%
\subsection{31.5 Measurable Performance
Targets}\label{measurable-performance-targets}}

Define Service Level Objectives (SLOs) in percentiles. * \textbf{p50
(Median)}: Typical case. * \textbf{p99}: The ``slow'' case (1 in 100). *
\textbf{p99.9}: The tail latency (1 in 1000). Crucial for HFT.

\textbf{Example Target}: ``Order processing must have p99 latency
\textless{} 5 microseconds.''

\begin{center}\rule{0.5\linewidth}{0.5pt}\end{center}
